% texte.tex — f170v, vv. 679–720 (conversion initiale des travaux antérieurs)
\vers[n=680, indent=2em, findent=1.2em]{}{}{μανίαις φλέγων λύσσῃ παράκο-}{celui d'eux qui sera pris à faire un sacrilège,}
\vers[n=681, indent=2em, findent=1.2em]{}{}{\alinea πος, εἴ τι δρῴη πᾶσιν ἐμφανὴς ὁρᾶν}{enflammé de fureur, délirant de rage,}
\vers[n=682, indent=2em, findent=1.2em]{}{}{ἔσται γυναιξὶ καὶ βροτοῖς,}{témoignera aux femmes et à tous les mortels}
\vers[n=683, indent=2em, findent=1.2em]{}{}{ὅτι τά παράνομα τά τ᾽ ἀνόσια παρὼν}{que la divinité châtie les atteintes aux lois}
\vers[n=685, indent=2em, findent=1.2em]{}{}{θεὸς παραχρῆμ᾽ ἀποτίνεται.}{et à la piété et ne se fait pas attendre.}
\vers[n=686, metre={tr⁴}]{ΚΟ.}{Coryphée}{ἀλλ᾽ ἔοιχ᾽ ἡμῖν ἅπαντά πως διεσκέφθαι καλῶς.}{Mais il nous semble avoir bien cherché partout et nous ne voyons pas d'autre homme caché parmi nous.}
\vers[n=687]{}{}{οὐχ ὁρῶμεν γοῦν ἔτ᾽ ἄλλον οὐδέν᾽ ἐγκαθήμενον}{\didasc{Mnésiloque enlève l'enfant et se réfugie sur l'autel.}}
\vers[n=688]{ΜΙ.}{Mika}{ἆ ἆ}{Ah ! Ah !}
\vers[suite]{ΜΙ.}{Mika}{ποῖ ποῖ σὺ φεύγεις; οὗτος οὗτος οὐ μενεῖς;}{Où ? Où te sauves-tu ? holà ! holà ! Arrête.}
\vers[n=690]{}{}{τάλαιν᾽ ἐγὼ τάλαινα, καὶ τὸ παιδίον}{Malheur ! Ah ! malheur ! le voilà parti avec mon enfant}
\vers[n=691]{}{}{ἐξαρπάσας μοι φροῦδος ἀπὸ τοῦ τιτθίου.}{arraché de ma mamelle.}
\vers[n=692]{ΜΝ.}{Mnésiloque}{κέκραχθι· τοῦτο δ᾽ οὐδέποτε σὺ ψωμιεῖς,}{Hurle. Jamais plus tu ne l'appateras,}
\vers[n=693]{}{}{ἢν μή μ᾽ ἀφῆτ᾽· ἀλλ᾽ ἐνθάδ᾽ ἐπὶ τῶν μηρίων}{si vous ne me laissez pas libre. Mais ici même « sur ces oblations,}
\vers[n=694]{}{}{πληγὲν μαχαίρᾳ τῇδε φοινίας φλέβας}{frappé de ce couteau, du pur sang de ses veines}
\vers[n=695]{}{}{καθαιματώσει βωμόν.}{il rougira l'autel ».}
\vers[suite]{ΜΙ.}{Mika}{ὢ τάλαιν᾽ ἐγώ.}{Malheur à moi !}
\vers[n=696]{}{}{γυναῖκες, οὐκ ἀρήξετ᾽; οὐ πολλὴν βοὴν}{Femmes, au secours ! Poussez une puissante clameur,}
\vers[n=697]{}{}{στήσεσθε καὶ τροπαῖον, ἀλλὰ τοῦ μόνου}{une clameur à faire lâcher pied. Souffrirez-vous}
\vers[n=698]{}{}{τέκνου με περιόψεσθ᾽ ἀποστερουμένην;}{qu'on me ravisse mon unique enfant?}
\vers[n=699, indent=2em, findent=1.2em, metre={lyr.}]{ΧΟ.}{Chœur}{ἔα ἔα.}{Ah! Ah !}
\vers[n=700, indent=2em, findent=1.2em]{}{}{ὦ πότνιαι Μοῖραι τί τόδε δέρκομαι}{Augustes M\oe res, quel prodige}
\vers[n=701, indent=2em, findent=1.2em]{}{}{νεοχμὸν αὖ τέρας;}{nouveau vient frapper mes yeux !}
\vers[n=702, metre={tr⁴}]{}{}{ὡς ἅπαντ᾽ ἄρ᾽ ἐστὶ τόλμης †πλέα† κἀναισχυντίας.}{Tout n'est donc qu'audace doublée d'effronterie !}
\vers[n=703]{}{}{οἷον αὖ δέδρακεν ἔργον, οἷον αὖ φίλαι τόδε.}{Quel forfait encore, amies, quel forfait il a commis !}
\vers[n=704]{ΜΝ.}{Mnésiloque}{οἷον; ὑμῶν ἐξαρήξω τὴν ἄγαν αὐθαδίαν.}{Quel forfait ? je vous extirperai votre outrecuidance !}
\vers[n=705]{ΧΟ.}{Chœur}{ταῦτα δῆτ᾽ οὐ δεινὰ πράγματ᾽ ἐστὶ καὶ περαιτέρω;}{N'est-ce pas une abomination, et pis encore ?}
\vers[n=706]{ΜΙ.}{Mika}{δεινὰ δῆθ᾽, ὁτιή γ᾽ ἔχει μου 'ξαρπάσας τὸ παιδίον.}{Une abomination, oui, qu'on m'ait arraché mon enfant.}
\vers[n=707, indent=3em, findent=1.7999999999999998em, metre={an.}]{ΧΟ.}{Chœur}{τί ἂν οὖν εἴποι πρὸς ταῦτά τις, ὅτε}{Que dire après cela, quand}
\vers[n=708, indent=3em, findent=1.7999999999999998em]{}{}{τοιαῦτα ποιῶν ὅδ᾽ ἀναισχυντεῖ;}{l'auteur d'une telle action a perdu toute honte ?}
\vers[n=709, indent=3em, findent=1.7999999999999998em]{ΜΝ.}{Mnésiloque}{κοὔπω μέντοι γε πέπαυμαι.}{Et je n'ai pas fini.}
\vers[n=710, indent=2em, findent=1.2em, metre={lyr.}]{ΧΟ.}{Chœur}{ἀλλ᾽ οὖν ἥκειν γ᾽ ὅθεν ἥκεις.}{Au moins, reviens d'où tu es parti ; ce n'est pas}
\vers[n=711, indent=2em, findent=1.2em]{}{}{φαύλως δ᾿ ἀποδρὰς οὐκ ἔξει.}{avec une fuite facile que tu t'en sortiras.}
\vers[n=712, indent=2em, findent=1.2em]{}{}{οἷον δράσας διέδυς ἔργον·}{Avec tout ce que tu as fait en t'immisçant ici,}
\vers[n=713, indent=2em, findent=1.2em]{}{}{λήψει δὲ κακόν.}{malheur t'en prendra !}
\vers[n=714, metre={tr⁴}]{ΜΝ.}{Mnésiloque}{τοῦτο μέντοι μὴ γένοιτο μηδαμῶς, ἀπεύχομαι.}{Puisse-t-il n'en être rien, j'en conjure le ciel.}
\vers[n=715, indent=2.5em, findent=1.5em, metre={do.}]{ΧΟ.}{Chœur}{τίς οὖν σοι, τίς ἂν σύμμαχος ἐκ θεῶν}{Et qui donc, qui des dieux immortels s'associera}
\vers[n=716, indent=2.5em, findent=1.5em]{}{}{ἀθανάτων ἔλθοι ξὺν ἀδίκοις ἔργοις;}{à des actes injurieux ?}
\vers[n=717]{ΜΝ.}{Mnésiloque}{μάτην λαλεῖτε· τὴνδ᾽ ἐγὼ οὐκ ἀφήσω.}{Inutiles bavardages ! Je ne la lâcherai pas.}
\vers[n=719, indent=2em, findent=1.2em, metre={lyr.}]{ΧΟ.}{Chœur}{ἀλλ᾽ οὐ μὰ τὼ θεὼ τάχ᾽ οὐ}{Ah! par les deux déesses, peut-être}
\vers[n=720, indent=2em, findent=1.2em]{}{}{χαίρων ἴσως ἐνυβριεῖς}{n'auras-tu pas lieu de rire de tes outrages}
\vers[indent=2em, findent=1.2em]{}{}{λόγους τε λέξεις ἀνοσίους}{et des discours insensés que tu profères.}
